\documentclass[preprint,numbers,1p]{elsarticle}

\usepackage{amsmath}
\usepackage{amssymb}
\usepackage{graphicx}
\usepackage{hyperref}
\usepackage{xcolor}
\usepackage{booktabs}
\usepackage{siunitx}
\usepackage{listings}
\usepackage{geometry}

\geometry{margin=2.5cm}

\lstset{
  basicstyle=\ttfamily\small,
  breaklines=true,
  frame=single,
  language=C++,
  keywordstyle=\color{blue},
  commentstyle=\color{gray},
}


\begin{document}

\begin{frontmatter}

\title{Physics Model of the DC-RSD Surface Simulator (\texttt{DCRSD\_SurfaceSimulator.c})}

\author{N. Cartiglia, INFN, Torino}

\begin{abstract}
This note documents the physics model implemented in \texttt{DCRSD\_SurfaceSimulator.c},
a ROOT-based GUI simulation for DC-coupled Resistive Silicon Detectors (RSD) with a $4\times4$ pixel
layout and optional resistive strips. The simulation solves the Kirchhoff current law (KCL)
on a sparse resistive network, models charge sharing via pixel templates, reconstructs
particle positions with a $\chi^2$ minimiser with Student's $t$ and Gaussian resolution fits,
provides inter-electrode resistance measurements through a Kelvin two-terminal method,
quantifies signal leakage to outer pixels, and supports both fixed laser charge and
MIP Landau charge fluctuation modes. All model parameters are controlled interactively
through a ROOT GUI.
\end{abstract}

\begin{keyword}
Resistive Silicon Detector \sep Resistive sheet \sep Charge sharing \sep Position resolution \sep ROOT simulation
\end{keyword}

\end{frontmatter}

%--------------------------------------------------------------------
\section{Overview}
\label{sec:overview}
%--------------------------------------------------------------------

The simulation models a $4\times4$ array of RSD pixels as a two-dimensional resistive network
on a uniform rectangular grid. Each grid node represents a small patch of the detector surface.
An ionising particle deposits charge at a given position; the charge then spreads through the
resistive layer and is collected by the pixel electrodes. The collected signals are used to
reconstruct the hit position.

The main simulation parameters are listed in Table~\ref{tab:params}.

\begin{table}[h]
\centering
\caption{GUI-controllable simulation parameters.}
\label{tab:params}
\begin{tabular}{lllr}
\toprule
Parameter & Symbol & Unit & Default \\
\midrule
Pixel pitch & $p$ & \si{\micro\meter} & 500 \\
Surface resistivity & $\rho_s$ & $\Omega/\square$ & 1000 \\
Strip resistivity at pad & $\rho_{s,\mathrm{pad}}$ & $\Omega/\square$ & 25 \\
Strip resistivity at mid & $\rho_{s,\mathrm{mid}}$ & $\Omega/\square$ & 25 \\
Strip gap at midpoint & $g_\mathrm{mid}$ & \si{\micro\meter} & 0 \\
Strip gap at pad & $g_\mathrm{pad}$ & \si{\micro\meter} & 0 \\
Pad radius & $a$ & \si{\micro\meter} & 60 \\
Input impedance & $R_\mathrm{in}$ & \si{\ohm} & 50 \\
Noise RMS & $\sigma_n$ & (normalised) & 0.018 \\
Scan steps per pixel & $n_s$ & -- & 100 \\
Events for template reco & $N_\mathrm{ev}$ & -- & 5000 \\
Trench between pixels & -- & boolean & false \\
No structure (surface only) & -- & boolean & false \\
MIP Landau charge & -- & boolean & false \\
\bottomrule
\end{tabular}
\end{table}

%--------------------------------------------------------------------
\section{Grid and Coordinate System}
\label{sec:grid}
%--------------------------------------------------------------------

The simulation domain covers a $3p \times 3p$ area centred on a $3\times3$ sub-array of the
$4\times4$ layout. The grid has $n = 301$ nodes along each axis, giving a node pitch (mesh
size) of
\begin{equation}
  h = \frac{3p}{n-1} = \frac{3p}{300}.
\end{equation}
For the default pitch $p = \SI{500}{\micro\meter}$ this gives $h = \SI{5}{\micro\meter}$ per
node. The pixel pitch occupies $p/h = 100$ nodes.

Node $(i,j)$ has physical coordinates
\begin{align}
  x_{ij} &= j \cdot h, & y_{ij} &= i \cdot h,
\end{align}
with $i,j \in \{0,\ldots,300\}$. The scalar node index is $\mathrm{idx} = i \cdot n + j$.
The total number of nodes is $N_\mathrm{tot} = n^2 = 90601$.

The four pixel electrodes in the central $2\times2$ sub-array are located at
$(r,c)\,p$ with $r,c\in\{1,2\}$ (using 0-based row/column index from 0 to 3). All 16 pad
positions are defined analogously.

%--------------------------------------------------------------------
\section{Resistive Network Model}
\label{sec:network}
%--------------------------------------------------------------------

\subsection{Edge conductance}
\label{sec:edge}

Each pair of adjacent nodes $(i,j)$ and $(i',j')$ is connected by an edge conductance
$g_{ij\to i'j'}$ that reflects the local material properties. Three contributions are
distinguished:

\paragraph{Surface resistive layer}
The resistive surface layer has a uniform surface resistivity $\rho_s$. For a square of side $h$
the surface conductance is
\begin{equation}
  g_\mathrm{surf} = \frac{1}{R_\mathrm{surf}}, \qquad
  R_\mathrm{surf} = \rho_s.
\end{equation}
(One square contributes $\rho_s$ regardless of $h$.)

\paragraph{Resistive strip}
Along certain inter-electrode paths a metallic strip with lower resistivity can be present.
The strip occupies a band of width $w_\mathrm{strip} = 2\times0.025\,p = 0.05\,p$ centred on
the electrode row or column. Within this band the strip conductance is added in parallel to the
surface:
\begin{equation}
  g = g_\mathrm{surf} + g_\mathrm{strip}, \qquad
  g_\mathrm{strip} = \frac{1}{r_\mathrm{strip}(d)},
\end{equation}
where $d$ is the distance of the edge midpoint from the nearest pad, and $r_\mathrm{strip}$
is the local strip sheet resistance interpolated linearly between its value at the pad
$(\rho_{s,\mathrm{pad}})$ and its value at the mid-point $(\rho_{s,\mathrm{mid}})$:
\begin{equation}
  r_\mathrm{strip}(d) = \rho_{s,\mathrm{pad}} + (\rho_{s,\mathrm{mid}} - \rho_{s,\mathrm{pad}})
  \frac{2d}{p}.
\end{equation}

Two gap regions suppress the strip contribution (returning only surface conductance):
\begin{itemize}
  \item \textbf{Midpoint gap} $g_\mathrm{mid}$: edges within half-width $g_\mathrm{mid}/2$ of
        the inter-electrode midpoint ($d = p/2$) carry only surface conductance.
  \item \textbf{Pad gap} $g_\mathrm{pad}$: edges within distance $g_\mathrm{pad}$ of any
        electrode ($d < g_\mathrm{pad}$) carry only surface conductance. This models a strip
        that ends before reaching the electrode, decoupling the strip from the electrode and
        increasing the effective inter-electrode resistance without degrading position resolution.
\end{itemize}

The number of squares from pad to pad along the strip (with no gaps) is
\begin{equation}
  N_\square = \frac{p}{w_\mathrm{strip}} = \frac{p}{0.05\,p} = 20.
\end{equation}
The total resistance from pad to pad is
\begin{equation}
  R_\mathrm{strip,total} = \frac{N_\square}{2}(\rho_{s,\mathrm{pad}} + \rho_{s,\mathrm{mid}})
  = 10\,(\rho_{s,\mathrm{pad}} + \rho_{s,\mathrm{mid}}).
\end{equation}
Due to the two-dimensional nature of the resistive network a correction factor of approximately
$1.7\times$ is observed with respect to the naive 1D (five-lane parallel) estimate.

\paragraph{Trench}
When the trench option is active, all strip-related conductance is suppressed and edges crossing
the inter-electrode boundaries are cut (zero conductance), except within a small radius around
each electrode. Both cardinal edges (purely horizontal or purely vertical) and diagonal edges
that cross a boundary are blocked; diagonal edges were historically omitted and represented a
source of small trench leakage. Only the surface contribution remains on non-trench edges. This
models an etched gap that completely interrupts current flow between pixels.

\paragraph{No structure (surface only)}
When the ``No structure'' option is active, the strip gap at midpoint is set equal to the full
pixel pitch $p$. Because the gap half-width $g_\mathrm{mid}/2 = p/2$ covers every inter-electrode
position, the strip contribution is suppressed everywhere and the surface behaves as a pure
uniform resistive surface. The strip resistance entries are disabled in the GUI to prevent
accidental modification. This mode is equivalent to a detector with no resistive strips and no
trench.

\subsection{KCL matrix assembly}
\label{sec:kcl}

The resistive network is described by the Kirchhoff Current Law (KCL) at every node. Denoting
the node voltage vector as $\mathbf{V}$ and the injected current vector as $\mathbf{B}$, the
system to be solved is
\begin{equation}
  \mathbf{M}\,\mathbf{V} = \mathbf{B},
\end{equation}
where $\mathbf{M}$ is the $N_\mathrm{tot}\times N_\mathrm{tot}$ sparse conductance (Laplacian)
matrix. Its off-diagonal entries are
\begin{equation}
  M_{\mathrm{idx},\,\mathrm{idx}'} = -g_{\mathrm{idx}\to\mathrm{idx}'},
\end{equation}
and the diagonal entries are the sum of all conductances leaving node $\mathrm{idx}$:
\begin{equation}
  M_{\mathrm{idx},\mathrm{idx}} = \sum_{\mathrm{neighbours}\,k} g_{\mathrm{idx}\to k}
  + \Delta_\mathrm{diag},
\end{equation}
where $\Delta_\mathrm{diag}$ is an additional diagonal term used to enforce boundary conditions
(see Section~\ref{sec:bc}).

\subsection{Boundary conditions at pads}
\label{sec:bc}

Two operating modes are available, selected by the input impedance $R_\mathrm{in}$:

\paragraph{$R_\mathrm{in} = 0$ (direct ground).}
Each pad node is forced to $V = 0$ by setting its diagonal entry to unity and leaving all
off-diagonal entries in that row at zero:
\begin{equation}
  M_{\mathrm{pad},\mathrm{pad}} = 1, \qquad
  M_{\mathrm{pad},k} = 0 \; \forall\, k\neq \mathrm{pad}, \qquad
  B_\mathrm{pad} = 0.
\end{equation}
This is equivalent to a short circuit to ground without requiring a large penalty conductance.

\paragraph{$R_\mathrm{in} > 0$ (floating pad with input impedance).}
The pad is not pinned; instead an additional conductance $1/R_\mathrm{in}$ is added between
the pad node and ground:
\begin{equation}
  M_{\mathrm{pad},\mathrm{pad}} \mathrel{+}= \frac{1}{R_\mathrm{in}}, \qquad
  B_\mathrm{pad} = 0.
\end{equation}
The collected current is then $I_k = V_k / R_\mathrm{in}$.

A high $R_\mathrm{in}$ isolates the readout electronics from the resistive strip while
preserving the charge-sharing structure needed for position reconstruction. The
inter-electrode resistance (Section~\ref{sec:resistance}) is independent of $R_\mathrm{in}$,
as it is measured with all pads floating.

\subsection{Matrix decomposition and caching}
\label{sec:cache}

The KCL matrix depends only on the geometric and material parameters (pitch, resistivities,
gaps, trench, $R_\mathrm{in}$) and not on the injection position. It is therefore factorised
once and reused for all injection positions. The decomposition uses \texttt{TDecompSparse}
(ROOT's sparse LU factoriser).

To avoid repeated decomposition between sessions, the matrix is cached on disk in a
sub-directory \texttt{rsd\_cache/}. Two cache key formats are used:
\begin{itemize}
  \item Without trench: \texttt{PP\{p\}\_Rsh\{$\rho_s$\}\_RsP\{$\rho_{s,P}$\}\_RsM\{$\rho_{s,M}$\}\_gap\{$g_\mathrm{mid}$\}\_gapP\{$g_\mathrm{pad}$\}\_tr0\_Rin\{$R_\mathrm{in}$\}\_Rpad\{$a$\}}
  \item With trench: \texttt{PP\{p\}\_Rsh\{$\rho_s$\}\_tr1\_Rin\{$R_\mathrm{in}$\}\_Rpad\{$a$\}} (strip parameters omitted as they have no effect)
\end{itemize}
On loading, the matrix is restored using the \texttt{TMatrixDSparse} copy constructor to
avoid ROOT's incompatibility check in the assignment operator.

%--------------------------------------------------------------------
\section{Signal Generation and Collection}
\label{sec:signal}
%--------------------------------------------------------------------

\subsection{Charge injection and fluctuation model}
\label{sec:injection}

A particle crossing the detector at position $(x_0, y_0)$ is modelled as a point current source
injected at the single nearest grid node, found by rounding to the closest integer node index:
\begin{equation}
  B_\mathrm{idx} = L \cdot I_\mathrm{total}, \qquad
  \mathrm{idx} = \mathrm{round}(x_0/h)\cdot n + \mathrm{round}(y_0/h),
\end{equation}
where $I_\mathrm{total} = 1$ (normalised) and $L$ is the amplitude factor described below.

\paragraph{Charge fluctuation modes.}
The GUI checkbox \emph{MIP Landau charge (laser = off)} controls the amplitude factor $L$:

\begin{itemize}
  \item \textbf{Laser mode} ($L=1$): the injected charge is fixed, corresponding to a laser
        pulse of constant intensity. All amplitude variation comes from electronics noise only.
  \item \textbf{MIP Landau mode}: $L$ is drawn from a Landau distribution with most-probable
        value $\mathrm{MPV}=1$ and scale $\sigma_L = 0.3$, independently for each simulated
        event:
        \begin{equation}
          L \sim \mathrm{Landau}(\mathrm{MPV}=1,\,\sigma_L=0.3).
        \end{equation}
        The noise is added to the signal \emph{after} scaling:
        \begin{equation}
          \hat{f}_k = L\cdot f_k(x,y) + \mathcal{N}(0,\sigma_n).
        \end{equation}
        The Landau fluctuations produce events with very small signal amplitude ($L \ll 1$),
        where the noise-to-signal ratio is large and position reconstruction fails.
        This results in heavy-tailed residual distributions qualitatively similar to those
        observed with MIP particles in test-beam data.
\end{itemize}

\subsection{Current collection}
\label{sec:collection}

After solving $\mathbf{M}\,\mathbf{V} = \mathbf{B}$, the current collected by pad $k$ is
computed as the net outflow from the pad node. Two formulas apply:

\begin{itemize}
  \item $R_\mathrm{in} = 0$: $I_k = \sum_{\mathrm{neighbours}\,m} g_{k\to m}(V_k - V_m)$
        (KCL edge sum).
  \item $R_\mathrm{in} > 0$: $I_k = V_k / R_\mathrm{in}$.
\end{itemize}

The fractional sharing for pixel $k$ is $f_k = I_k / \sum_j I_j$.

%--------------------------------------------------------------------
\section{Transient Signal Propagation}
\label{sec:transient}
%--------------------------------------------------------------------

An optional transient model propagates the signal in time. Each node has a capacitance
\begin{equation}
  C_\mathrm{node} = \frac{\varepsilon_0\,\varepsilon_\mathrm{Si}\,h^2}{d_\mathrm{Si}},
\end{equation}
where $d_\mathrm{Si}$ is the silicon thickness and $h$ is the node pitch. The time step is
$\Delta t = \SI{20}{\pico\second}$. At each time step the node voltages evolve according to the
RC equation:
\begin{equation}
  C_\mathrm{node}\,\frac{dV_k}{dt} = -\sum_m g_{k\to m}(V_k - V_m) + B_k(t).
\end{equation}
This is integrated with a forward-Euler scheme.

Signal propagation delays are also computed geometrically using a Dijkstra--Fermat
shortest-path algorithm on the resistive network, giving the time for a wavefront to travel
from the injection point to each pixel electrode.

%--------------------------------------------------------------------
\section{Validity of the Quasi-Static Approximation}
\label{sec:quasistatic}
%--------------------------------------------------------------------

The simulation solves a static (DC) resistive network. Real detector signals from
LGAD/RSD devices are short current pulses with significant spectral content between
approximately \SI{50}{MHz} and \SI{400}{MHz}. This section quantifies the error
introduced by neglecting the finite-frequency response of the resistive sheet.

\subsection{Bulk capacitive shunting}
\label{sec:bulk_cap}

At frequency $\omega = 2\pi f$, each node of the resistive sheet is capacitively
coupled to the back-plane electrode through the silicon bulk. The capacitance per unit area is
\begin{equation}
  C_\mathrm{bulk} = \frac{\varepsilon_0\,\varepsilon_\mathrm{Si}}{d_\mathrm{Si}},
\end{equation}
where $\varepsilon_\mathrm{Si} = 11.7$ and $d_\mathrm{Si}$ is the detector thickness.
This shunt admittance modifies the 2D resistive diffusion equation into a
Helmholtz-type equation for the sheet potential:
\begin{equation}
  \nabla^2 V - \frac{j\omega\,C_\mathrm{bulk}}{\rho_s^{-1}}\,V = -\rho_s\,j_\mathrm{inj},
\end{equation}
where $j_\mathrm{inj}$ is the injected current density.
Solutions decay with a complex \emph{AC spreading length}
\begin{equation}
  \lambda_\mathrm{AC} = \frac{1}{\sqrt{j\omega\,C_\mathrm{bulk}\,\rho_s}},
  \qquad
  |\lambda_\mathrm{AC}| = \frac{1}{\sqrt{\omega\,C_\mathrm{bulk}\,\rho_s}}.
\end{equation}
The DC limit corresponds to $|\lambda_\mathrm{AC}| \to \infty$. The fractional correction
to the signal sharing at distance $p$ from the injection point is of order
$(p/|\lambda_\mathrm{AC}|)^2$.

\subsection{Numerical estimates}
\label{sec:quasistatic_numbers}

For a typical DC-RSD detector with $d_\mathrm{Si} = \SI{300}{\micro\meter}$ and
$\rho_s = \SI{1000}{\ohm/\square}$:
\begin{equation}
  C_\mathrm{bulk} = \frac{8.85\times10^{-12}\times 11.7}{300\times10^{-6}}
  \approx \SI{3.45e-8}{F/m^2}.
\end{equation}
Table~\ref{tab:lambda} lists $|\lambda_\mathrm{AC}|$ and the corresponding sharing
correction for a pixel pitch $p = \SI{450}{\micro\meter}$.

\begin{table}[h]
\centering
\caption{AC spreading length and fractional correction to signal sharing
  for $\rho_s = \SI{1000}{\ohm/\square}$, $d_\mathrm{Si} = \SI{300}{\micro\meter}$,
  $p = \SI{450}{\micro\meter}$.}
\label{tab:lambda}
\begin{tabular}{rrrr}
\toprule
$f$ & $|\lambda_\mathrm{AC}|$ & $p/|\lambda_\mathrm{AC}|$ & Correction $(p/\lambda)^2$ \\
\midrule
\SI{50}{MHz}  & \SI{9.6}{mm}  & 0.047 & $\sim 0.2\%$ \\
\SI{100}{MHz} & \SI{6.8}{mm}  & 0.066 & $\sim 0.4\%$ \\
\SI{200}{MHz} & \SI{4.8}{mm}  & 0.094 & $\sim 0.9\%$ \\
\SI{400}{MHz} & \SI{3.4}{mm}  & 0.13  & $\sim 1.7\%$ \\
\bottomrule
\end{tabular}
\end{table}

Since $|\lambda_\mathrm{AC}| \gg p$ across the full signal bandwidth, the spatial
distribution of collected charge is virtually indistinguishable from the DC case.

\subsection{Strip RC time constant}
\label{sec:strip_rc}

The resistive strips (Section~\ref{sec:edge}) have a typical resistance
$R_\mathrm{strip} \sim 500\,\Omega$ for the pad-to-pad path. Their self-capacitance
to the back-plane is negligible ($C_\mathrm{strip} \sim \SI{0.1}{fF}$), giving a
characteristic frequency
\begin{equation}
  f_\mathrm{3dB} = \frac{1}{2\pi\,R_\mathrm{strip}\,C_\mathrm{strip}} \sim \SI{3}{THz},
\end{equation}
which is far above the signal bandwidth. The strip impedance is therefore purely
resistive at all relevant frequencies.

\subsection{Conclusion}
\label{sec:quasistatic_conclusion}

The quasi-static (DC) approximation is valid for DC-RSD signal sharing at frequencies
up to at least \SI{400}{MHz}. The maximum fractional error in signal sharing is
$\lesssim 2\%$ over the full LGAD/RSD signal bandwidth, well below the level of
noise and other systematic effects. No frequency-dependent correction to the
resistive network model is required under typical detector parameters.

%--------------------------------------------------------------------
\section{Charge Sharing Templates}
\label{sec:templates}
%--------------------------------------------------------------------

The sharing template describes how the collected charge fraction $f_k(x,y)$ varies as a
function of the hit position $(x,y)$ within the central pixel cell. Templates are computed on
a $(n_s+1)\times(n_s+1)$ grid spanning one pixel pitch in each direction, with step size
$\Delta x = p/n_s$. The GUI parameter $n_s$ (default 100) sets the number of scan steps per
pixel; at the default, $\Delta x = h$ (one grid node per step). For each scan point the KCL
system is solved and the fractions $\{f_k\}$ for all 16 pixels are stored. The resulting
$f_k(x,y)$ maps are displayed as 2D colour histograms (\texttt{TH2D}) with $z$ range $[0,1]$.

%--------------------------------------------------------------------
\section{Signal Leakage}
\label{sec:leakage}
%--------------------------------------------------------------------

For a hit inside the central pixel cell $[p,2p]\times[p,2p]$, the signal collected by the
four central pixels (P6, P7, P10, P11) should ideally equal the total injected signal. Any
current reaching the 12 outer pixels constitutes \emph{signal leakage}:
\begin{equation}
  \mathcal{L}(x,y) = 1 - \frac{I_6 + I_7 + I_{10} + I_{11}}{I_\mathrm{total}}.
\end{equation}
The leakage map is computed over the central pixel cell by scanning at intervals of $\Delta x$
and displayed alongside the central-pad fraction map. The surface-averaged leakage,
\begin{equation}
  \langle\mathcal{L}\rangle = \frac{1}{N_\mathrm{bins}}\sum_{\mathrm{bins}} \mathcal{L}(x,y),
\end{equation}
is printed to the terminal after the scan. The two maps share the same scan and are computed
simultaneously without additional solver calls.

Leakage is largest near the pixel boundaries, where the signal is shared between pixels, and
is enhanced when a resistive strip at low resistivity provides a direct current path to outer
electrodes. The combination of a high $R_\mathrm{in}$ with a pad gap $g_\mathrm{pad}$ reduces
leakage while preserving the position resolution.

%--------------------------------------------------------------------
\section{Position Reconstruction}
\label{sec:reco}
%--------------------------------------------------------------------

\subsection{Template-based $\chi^2$ minimisation}
\label{sec:chi2}

Given a measured set of signal fractions $\{\hat{f}_k\}$ (with Gaussian noise of RMS
$\sigma_n$ added), the reconstructed position $(\hat{x}, \hat{y})$ is found by minimising
\begin{equation}
  \chi^2(x,y) = \sum_{k \in \{6,7,10,11\}} [\hat{f}_k - f_k(x,y)]^2,
\end{equation}
where the sum runs over the \textbf{four central pads} only (P6, P7, P10, P11), since the
templates are built exclusively for these pixels.
The minimisation scans the $(n_s+1)\times(n_s+1)$ template grid to find the coarse minimum, then applies
parabolic interpolation in both $x$ and $y$ independently to obtain sub-grid precision:
\begin{equation}
  x^* = x_m - \frac{1}{2}\,\frac{\chi^2_{m+1} - \chi^2_{m-1}}{\chi^2_{m+1} - 2\chi^2_m + \chi^2_{m-1}}\,\Delta x.
\end{equation}
Parabolic interpolation removes the scan-direction bias that would otherwise produce an
asymmetric $x$-residual distribution.

\subsection{Noiseless reconstruction and distortion grid}
\label{sec:noiseless}

Before the noisy event loop, a noiseless reconstructed position is computed analytically:
\begin{align}
  x_\mathrm{nl} &= p\cdot\frac{3}{2} + \frac{(f_{10}+f_{11})-(f_6+f_7)}{f_6+f_7+f_{10}+f_{11}}\,\frac{p}{2}, \\
  y_\mathrm{nl} &= p\cdot\frac{3}{2} + \frac{(f_7+f_{11})-(f_6+f_{10})}{f_6+f_7+f_{10}+f_{11}}\,\frac{p}{2},
\end{align}
using the unbalance formula with noiseless amplitudes. The noiseless reconstruction is used
to populate the \emph{distortion grid} ($x_\mathrm{true}$ vs $x_\mathrm{nl}$) as a scatter
plot, which shows the systematic nonlinearity of the unbalance estimator without contamination
from noise. A separate noisy realisation is used to build the distortion map displayed as a
2D histogram.

\subsection{Resolution map}
\label{sec:resmap}

The position resolution is evaluated by generating $N_\mathrm{ev}$ random hit positions uniform
in $[p, 2p] \times [p, 2p]$ (the central pixel cell). For each hit the reconstruction error is
$\epsilon_x = \hat{x} - x_\mathrm{true}$. The RMS resolution map is computed on a
$(n_s+1)\times(n_s+1)$ grid, where each bin accumulates the squared errors:
\begin{equation}
  \sigma_x(x,y) = \sqrt{\frac{1}{N_\mathrm{bin}}\sum_{i\in\mathrm{bin}} \epsilon_{x,i}^2}.
\end{equation}
This definition includes both the statistical spread and any residual bias; it is equivalent
to the root-mean-square position error. The $z$ scale is fixed at $[0, \SI{20}{\micro\meter}]$.

\subsection{Resolution fits}
\label{sec:fits}

The global $x$-residual distribution $\Delta x = \hat{x} - x_\mathrm{true}$ is fitted with
two models displayed in separate panels:

\paragraph{Gaussian fit.}
A standard Gaussian $\mathcal{N}(\mu,\sigma)$ is fitted to extract the core resolution
$\sigma$.

\paragraph{Student's $t$ fit.}
A Student's $t$ distribution with location $\mu$, scale $\sigma_\mathrm{scale}$, and
$\nu$ degrees of freedom is fitted over the range $[-\SI{100}{\micro\meter},\,+\SI{100}{\micro\meter}]$:
\begin{equation}
  f(x;\,\mu,\sigma_\mathrm{scale},\nu) =
  \frac{1}{\sigma_\mathrm{scale}}\,
  t\!\left(\frac{x-\mu}{\sigma_\mathrm{scale}};\,\nu\right),
\end{equation}
where $t(z;\nu)$ is the standard Student's $t$ density implemented via
\texttt{TMath::Student}. For $\nu > 2$ the true standard deviation of the distribution is
\begin{equation}
  \sigma_\mathrm{RMS} = \sigma_\mathrm{scale}\sqrt{\frac{\nu}{\nu-2}}.
\end{equation}
The $t$ distribution accounts for heavy tails arising from ambiguous template matches near
pixel boundaries and from Landau amplitude fluctuations (MIP mode). As $\nu\to\infty$ it
converges to the Gaussian. Fit parameters
($\sigma_\mathrm{scale}$, $\nu$, $\sigma_\mathrm{RMS}$) are reported on the terminal.

\subsection{Canvas layouts}
\label{sec:linearity}

\paragraph{Unbalance analysis (canvas c3, $2\times2$).}
\begin{enumerate}
  \item[P1] \textbf{Distortion grid}: scatter of noiseless reconstructed positions
        $(x_\mathrm{nl}, y_\mathrm{nl})$ over the central pixel cell, showing the systematic
        nonlinearity of the unbalance estimator.
  \item[P2] \textbf{2D RMS resolution map}: $\sigma_x(x,y)$ on a $(n_s+1)\times(n_s+1)$ grid ($z$ range
        $[0, \SI{20}{\micro\meter}]$).
  \item[P3] \textbf{1D residuals} $\Delta x = \hat{x} - x_\mathrm{true}$, fitted with a
        Student's $t$ distribution. Fit parameters printed to the terminal.
  \item[P4] \textbf{X linearity scatter}: $x_\mathrm{true}$ vs.\ $x_\mathrm{nl}$ (noiseless),
        one point per scan grid position. Dashed red diagonal at $x_\mathrm{reco}=x_\mathrm{true}$.
\end{enumerate}

\paragraph{Template reconstruction analysis (canvas c5, $3\times2$).}
The canvas is divided into six panels ($3$ columns $\times$ $2$ rows):
\begin{enumerate}
  \item[P1] \textbf{Reconstructed position density}: 2D histogram of $(x_\mathrm{reco},
        y_\mathrm{reco})$ for all accepted events.
  \item[P2] \textbf{2D RMS resolution map}: same definition as in c3 ($z$ range
        $[0, \SI{30}{\micro\meter}]$).
  \item[P3] \textbf{RMS($\Delta x$) vs.\ $X_\mathrm{true}$}: 1D projection of the resolution
        map integrated over $Y$, with error bars $\sigma_\mathrm{RMS}/\sqrt{2N}$.
  \item[P4] \textbf{1D residuals} $\Delta x$ fitted with a Gaussian; $\sigma$ annotated on
        the pad.
  \item[P5] \textbf{1D residuals} (clone) fitted with a Student's $t$ distribution;
        $\sigma_\mathrm{scale}$, $\nu$, and $\sigma_\mathrm{RMS}$ reported on the terminal.
  \item[P6] \textbf{X linearity scatter}: $x_\mathrm{true}$ vs.\ $x_\mathrm{reco}$ (with
        noise), one point per event. Dashed red diagonal overlaid.
\end{enumerate}

Nonlinearity is an intrinsic feature of the unbalance estimator near pixel boundaries, where
the charge sharing changes rapidly. The linearity plot provides a direct visualisation of the
estimator's accuracy independent of the noise level.

%--------------------------------------------------------------------
\section{Inter-electrode Resistance Measurement}
\label{sec:resistance}
%--------------------------------------------------------------------

The effective resistance between two pixel electrodes is measured using a Kelvin (two-terminal)
method. For each pair of pads $(P_a, P_b)$ to be measured:

\begin{enumerate}
  \item Build the free KCL matrix with all pad nodes floating (no boundary conditions applied to pads). This matrix is independent of $R_\mathrm{in}$.
  \item Add a weak reference conductance $G_\mathrm{ref} = 10^{-9}$ to node 0 to remove the
        singular null mode:
        \begin{equation}
          M_{0,0} \mathrel{+}= G_\mathrm{ref}.
        \end{equation}
  \item Set the right-hand side $\mathbf{B}$ with $B(P_a) = +1$ (current source) and
        $B(P_b) = -1$ (current sink).
  \item Solve $\mathbf{M}\,\mathbf{V} = \mathbf{B}$.
  \item The inter-electrode resistance is $R_{ab} = V(P_a) - V(P_b)$.
\end{enumerate}

This approach leaves all other pads floating so that their presence does not short-circuit the
measurement path. Two measurements are performed:

\begin{itemize}
  \item $R_\mathrm{side}$: pixel $P_{11}$ (row 2, col 2) to $P_7$ (row 1, col 2) --- side neighbours.
  \item $R_\mathrm{diag}$: pixel $P_{11}$ (row 2, col 2) to $P_6$ (row 1, col 1) --- diagonal neighbours.
\end{itemize}

In a pure 2D resistive sheet the resistance between two electrodes scales logarithmically with
their separation, so the ratio $R_\mathrm{diag}/R_\mathrm{side}$ is
expected to be slightly above unity:
\begin{equation}
  \frac{R_\mathrm{diag}}{R_\mathrm{side}} \approx \frac{\ln(\sqrt{2}\,p/a)}{\ln(p/a)},
\end{equation}
where $a$ is the effective electrode radius. Numerically this gives a ratio of approximately
$1.13$--$1.15$, consistent with simulation results.

%--------------------------------------------------------------------
\section{Summary of Simulation Workflow}
\label{sec:workflow}
%--------------------------------------------------------------------

\begin{enumerate}
  \item \textbf{Grid setup}: generate the $301\times301$ node grid for the chosen pitch $p$.
  \item \textbf{Matrix build}: assemble the KCL conductance matrix \textbf{M} using
        \texttt{GetEdgeConductance()} for each of the four neighbour directions at every node.
        Apply pad boundary conditions according to $R_\mathrm{in}$.
  \item \textbf{Decomposition}: factorise \textbf{M} with \texttt{TDecompSparse}.
        Cache the matrix and decomposition on disk using a parameter-encoded file name.
  \item \textbf{Template scan}: for each of the $(n_s+1)\times(n_s+1)$ scan positions within the central
        pixel cell, inject current, solve, collect currents, store $\{f_k(x,y)\}$.
  \item \textbf{Signal leakage scan}: for each scan position in the central pixel cell, compute
        the fraction of signal reaching outer pixels; report the surface-averaged leakage.
  \item \textbf{Noiseless reconstruction}: for each scan point compute $x_\mathrm{nl}$, $y_\mathrm{nl}$
        using the unbalance formula with noiseless amplitudes; populate the linearity scatter plot.
  \item \textbf{Event loop}: generate $N_\mathrm{ev}$ random hits; optionally draw amplitude factor
        $L$ from a Landau distribution (MIP mode) or set $L=1$ (laser mode); add Gaussian noise;
        reconstruct positions via $\chi^2$ minimisation with parabolic interpolation.
  \item \textbf{Resolution map}: bin reconstruction errors into $(n_s+1)\times(n_s+1)$ map, compute
        RMS per bin.
  \item \textbf{Resolution fits}: fit the global $\Delta x$ distribution over $[-100,+100]\,\mu$m
        with a Gaussian and a Student's $t$ (via \texttt{TMath::Student});
        report $\sigma$, $\sigma_\mathrm{scale}$, $\nu$, $\sigma_\mathrm{RMS}$.
  \item \textbf{Linearity}: display $x_\mathrm{true}$ vs $x_\mathrm{reco}$ scatter with diagonal reference.
  \item \textbf{Resistance}: Kelvin measurement between selected pad pairs (independent of $R_\mathrm{in}$).
\end{enumerate}


%--------------------------------------------------------------------
\section*{Code and Documentation Availability}
\label{sec:availability}
%--------------------------------------------------------------------

The source code \texttt{DCRSD\_SurfaceSimulator.c} and this technical note are publicly
available on GitHub:
\begin{center}
  \url{https://github.com/cartiglia/DCRSD_SS}
\end{center}
To download the code, clone the repository with:
\begin{lstlisting}[language=bash]
git clone https://github.com/cartiglia/DCRSD_SS.git
\end{lstlisting}
The simulator requires ROOT 6 (\url{https://root.cern}) and is run directly as a ROOT macro:
\begin{lstlisting}[language=bash]
root DCRSD_SurfaceSimulator.c
\end{lstlisting}

\end{document}
